\section*{Annexe A -- Variables d'entrée du modèle}
\addcontentsline{toc}{section}{Annexe A -- Variables d'entrée du modèle}

Le tableau~\ref{tab:annexe-features} présente les principales constantes vitales et variables
utilisées en entrée du modèle de triage.

\begin{table}[h]
\centering
\renewcommand{\arraystretch}{1.3}
\begin{tabular}{p{3.2cm} p{3cm} p{7cm}}
\toprule
\textbf{Variable} & \textbf{Unité / Type} & \textbf{Description} \\
\midrule
\texttt{age}        & années & Âge du patient \\
\texttt{mental}     & catégoriel & État mental (alerte, voix, douleur, inconscient) \\
\texttt{nrsPain}    & 0--10 & Intensité de la douleur (échelle numérique NRS) \\
\texttt{sbp}        & mmHg & Pression artérielle systolique \\
\texttt{dbp}        & mmHg & Pression artérielle diastolique \\
\texttt{hr}         & bpm & Fréquence cardiaque \\
\texttt{rr}         & /min & Fréquence respiratoire \\
\texttt{bt}         & \textdegree C & Température corporelle \\
\texttt{saturation} & \% & Saturation pulsée en oxygène (SpO\textsubscript{2}) \\
\texttt{arrivalMode}& catégoriel & Mode d'arrivée aux urgences \\
\bottomrule
\end{tabular}
\caption{Principales variables d'entrée du modèle QuickTriage.}
\label{tab:annexe-features}
\end{table}

\section*{Annexe B -- Structure des métadonnées de prétraitement}
\addcontentsline{toc}{section}{Annexe B -- Structure des métadonnées de prétraitement}

Le fichier \texttt{metadata\_v3.json} externalise les paramètres de prétraitement chargés par le
backend au démarrage. Sa structure est la suivante :

\begin{listing}[H]
\begin{minted}[bgcolor=bg,fontsize=\small]{json}
{
  "model_version": "v3",
  "feature_order": ["age", "mental", "nrsPain", "sbp", "dbp",
                    "hr", "rr", "bt", "saturation", "arrivalMode", "..."],
  "imputer_medians": [/* mediane par feature */],
  "scaler_mean":     [/* moyenne par feature (StandardScaler) */],
  "scaler_scale":    [/* ecart-type par feature (StandardScaler) */],
  "seuil_critique":  0.30,
  "classes": ["CRITIQUE", "URGENT", "NON_URGENT"]
}
\end{minted}
\caption{Structure du fichier \texttt{metadata\_v3.json}.}
\label{lst:annexe-meta}
\end{listing}

\section*{Annexe C -- Correspondance KTAS / classes métier}
\addcontentsline{toc}{section}{Annexe C -- Correspondance KTAS / classes métier}

\begin{table}[h]
\centering
\renewcommand{\arraystretch}{1.4}
\begin{tabular}{p{3.2cm} p{2.5cm} p{6.5cm}}
\toprule
\textbf{Classe métier} & \textbf{Niveaux KTAS} & \textbf{Délai de prise en charge} \\
\midrule
\textcolor{qtcritique}{\textbf{CRITIQUE}}     & KTAS 1--2 & Immédiate ($<$ 15 min) \\
\textcolor{qturgent}{\textbf{URGENT}}         & KTAS 3    & Rapide \\
\textcolor{qtnonurgent}{\textbf{NON\_URGENT}} & KTAS 4--5 & Différée \\
\bottomrule
\end{tabular}
\caption{Correspondance entre les niveaux KTAS et les classes métier de QuickTriage.}
\label{tab:annexe-ktas}
\end{table}

\section*{Annexe D -- Hyperparamètres optimaux du modèle}
\addcontentsline{toc}{section}{Annexe D -- Hyperparamètres optimaux du modèle}

Le tableau~\ref{tab:annexe-hp} récapitule les hyperparamètres du meilleur essai Optuna
(essai n\textsuperscript{o}~64) retenus pour le modèle XGBoost final.

\begin{table}[h]
\centering
\renewcommand{\arraystretch}{1.3}
\begin{tabular}{p{5cm} c}
\toprule
\textbf{Hyperparamètre} & \textbf{Valeur} \\
\midrule
\texttt{n\_estimators}  & 458 \\
\texttt{max\_depth}     & 7 \\
\texttt{learning\_rate} & $\approx$ 0,0113 \\
Nombre d'essais Optuna  & 80 \\
Meilleur essai          & n\textsuperscript{o} 64 \\
\bottomrule
\end{tabular}
\caption{Hyperparamètres optimaux du modèle XGBoost final.}
\label{tab:annexe-hp}
\end{table}
